% --- Archon physics formalization source begin ---
% archon:physics
% archon:covers PhyXMiniProblems/problem_phyx_mini_0955.lean
% archon:source-report reports/phyx_mini/problem_phyx_mini_0955.source.json

\chapter{Physics problem phyx\_mini\_0955}
\label{ch:PhyXMiniProblems_problem_phyx_mini_0955}

\paragraph{Problem source.}
\#\# Physical scenario

You are a research scientist working on a high-energy particle accelerator. Using a modern version of the Thomson \textbackslash{}( e/m \textbackslash{}) apparatus, you want to measure the mass of a muon (a fundamental particle that has the same charge as an electron but greater mass). The magnetic field between the two charged plates is \textbackslash{}( 0.340\textbackslash{},\textbackslash{}text\{T\} \textbackslash{}). You measure the electric field for zero particle deflection as a function of the accelerating potential \textbackslash{}( V \textbackslash{}). This potential is large enough that you can assume the initial speed of the muons to be zero. An \textbackslash{}( E\textasciicircum{}2 \textbackslash{})-versus-\textbackslash{}( V \textbackslash{}) graph of your data.

\#\# Question

Use the graph to calculate the mass \textbackslash{}( m \textbackslash{}) of a muon.

\#\# Answer choices

- A: A: 1.63 \textbackslash{}times 10\textasciicircum{}\{-28\}\textbackslash{}

- B: B: 1.85 \textbackslash{}times 10\textasciicircum{}\{-28\}\textbackslash{}

- C: C: 2.25 \textbackslash{}times 10\textasciicircum{}\{-28\}\textbackslash{}

- D: D: 9.11 \textbackslash{}times 10\textasciicircum{}\{-31\}\textbackslash{}

\#\# Figure caption (auxiliary; use the image as primary evidence)

The image shows a graph with a linear relationship between two variables. The x-axis is labeled \textbackslash{}(V \textbackslash{}, (\textbackslash{}text\{V\})\textbackslash{}), representing voltage in volts, with intervals marked at 0, 100, 200, 300, and 400. The y-axis is labeled \textbackslash{}(E\textasciicircum{}2 \textbackslash{}, (10\textasciicircum{}8 \textbackslash{}, \textbackslash{}text\{V\}\textasciicircum{}2/\textbackslash{}text\{m\}\textasciicircum{}2)\textbackslash{}), representing the square of the electric field in units of \textbackslash{}(10\textasciicircum{}8 \textbackslash{}, \textbackslash{}text\{V\}\textasciicircum{}2/\textbackslash{}text\{m\}\textasciicircum{}2\textbackslash{}), with intervals at 0, 200, 400, 600, and 800.

The graph has a series of black dots representing data points, which form a straight line, indicating a strong linear relationship. A magenta line connects the dots, emphasizing the linearity. The relationship suggests that as voltage increases, \textbackslash{}(E\textasciicircum{}2\textbackslash{}) increases proportionally. The grid lines in the background help to align and determine the values of the data points.

\#\# Dataset metadata

Electromagnetic Induction; Multi-Formula Reasoning, Physical Model Grounding Reasoning

\paragraph{Recorded answer/context.}
B: B: 1.85 \textbackslash{}times 10\textasciicircum{}\{-28\}\textbackslash{}

\paragraph{Figure/image path.}
/root/proposal\_for\_physic/science-mango/phyx\_mini\_run/phyx\_data/test\_image/955.png

\paragraph{Formalization target.}
create a compiling Lean file with sorry bodies at `PhyXMiniProblems/problem\_phyx\_mini\_0955.lean`. The Lean declarations must preserve the physical quantities, dimensions or dimensional roles, figure labels, governing-law hypotheses, and final relation expressed by this problem.
Use Mathlib/Physlib names found through LeanExplore where available. If a physics API is missing, introduce faithful local abstractions rather than scalar placeholder aliases.

\begin{theorem}[Physics formalization target]
\label{thm:physics:phyx_mini_0955:target}
\lean{PhyXMiniProblems.ProblemPhyXMini0955.muonMass_agrees_with_choice_B}
\uses{def:physics:phyx-mini-0955:phyxminiproblems-problemphyxmini0955-massinkilograms, def:physics:phyx-mini-0955:phyxminiproblems-problemphyxmini0955-muonthomsonexperiment, def:physics:phyx-mini-0955:phyxminiproblems-problemphyxmini0955-matchesscenarioandgraphreadouts, def:physics:phyx-mini-0955:phyxminiproblems-problemphyxmini0955-haselementarychargecalibration, def:physics:phyx-mini-0955:phyxminiproblems-problemphyxmini0955-satisfiesmuonthomsonlaws, def:physics:phyx-mini-0955:phyxminiproblems-problemphyxmini0955-hasnondegeneratemuonrun, def:physics:phyx-mini-0955:phyxminiproblems-problemphyxmini0955-answerchoicemassinkilograms, def:physics:phyx-mini-0955:phyxminiproblems-problemphyxmini0955-agreestodisplayedprecision}
The assigned autoformalize agent should translate this physics problem into Lean declarations in the covered file, with theorem and lemma proof bodies written as `by sorry`.
\end{theorem}
\begin{proof}
This is an autoformalization task, not a proof task. Produce statements that can later be proved by the physics prover without weakening the physical model.
\end{proof}

% --- Archon declaration topology begin ---
\section{Declaration topology}
\begin{definition}[speed Dimension]
  \label{def:physics:phyx-mini-0955:phyxminiproblems-problemphyxmini0955-speeddimension}
  \lean{PhyXMiniProblems.ProblemPhyXMini0955.speedDimension}
  A speed has physical dimension L T⁻¹
\end{definition}
\begin{proof}
  This is the declared model definition of speed Dimension.
\end{proof}
\begin{definition}[electric Potential Dimension]
  \label{def:physics:phyx-mini-0955:phyxminiproblems-problemphyxmini0955-electricpotentialdimension}
  \lean{PhyXMiniProblems.ProblemPhyXMini0955.electricPotentialDimension}
  Electric potential has physical dimension M L² T⁻² C⁻¹
\end{definition}
\begin{proof}
  This is the declared model definition of electric Potential Dimension.
\end{proof}
\begin{definition}[electric Field Dimension]
  \label{def:physics:phyx-mini-0955:phyxminiproblems-problemphyxmini0955-electricfielddimension}
  \lean{PhyXMiniProblems.ProblemPhyXMini0955.electricFieldDimension}
  \uses{def:physics:phyx-mini-0955:phyxminiproblems-problemphyxmini0955-electricpotentialdimension}
  Electric-field magnitude has dimension M L T⁻² C⁻¹ (volts per metre)
\end{definition}
\begin{proof}
  This is the declared model definition of electric Field Dimension.
\end{proof}
\begin{definition}[magnetic Flux Density Dimension]
  \label{def:physics:phyx-mini-0955:phyxminiproblems-problemphyxmini0955-magneticfluxdensitydimension}
  \lean{PhyXMiniProblems.ProblemPhyXMini0955.magneticFluxDensityDimension}
  Magnetic flux density has dimension M T⁻¹ C⁻¹ (tesla)
\end{definition}
\begin{proof}
  This is the declared model definition of magnetic Flux Density Dimension.
\end{proof}
\begin{definition}[electric Field Squared Dimension]
  \label{def:physics:phyx-mini-0955:phyxminiproblems-problemphyxmini0955-electricfieldsquareddimension}
  \lean{PhyXMiniProblems.ProblemPhyXMini0955.electricFieldSquaredDimension}
  \uses{def:physics:phyx-mini-0955:phyxminiproblems-problemphyxmini0955-electricfielddimension}
  The vertical graph variable has the squared electric-field dimension
\end{definition}
\begin{proof}
  This is the declared model definition of electric Field Squared Dimension.
\end{proof}
\begin{definition}[Mass Quantity]
  \label{def:physics:phyx-mini-0955:phyxminiproblems-problemphyxmini0955-massquantity}
  \lean{PhyXMiniProblems.ProblemPhyXMini0955.MassQuantity}
  A nonnegative, unit-independent particle mass
\end{definition}
\begin{proof}
  This is the declared model definition of Mass Quantity.
\end{proof}
\begin{definition}[Charge Magnitude Quantity]
  \label{def:physics:phyx-mini-0955:phyxminiproblems-problemphyxmini0955-chargemagnitudequantity}
  \lean{PhyXMiniProblems.ProblemPhyXMini0955.ChargeMagnitudeQuantity}
  A nonnegative, unit-independent electric-charge magnitude
\end{definition}
\begin{proof}
  This is the declared model definition of Charge Magnitude Quantity.
\end{proof}
\begin{definition}[Speed Magnitude Quantity]
  \label{def:physics:phyx-mini-0955:phyxminiproblems-problemphyxmini0955-speedmagnitudequantity}
  \lean{PhyXMiniProblems.ProblemPhyXMini0955.SpeedMagnitudeQuantity}
  \uses{def:physics:phyx-mini-0955:phyxminiproblems-problemphyxmini0955-speeddimension}
  A nonnegative, unit-independent speed magnitude
\end{definition}
\begin{proof}
  This is the declared model definition of Speed Magnitude Quantity.
\end{proof}
\begin{definition}[Electric Potential Magnitude Quantity]
  \label{def:physics:phyx-mini-0955:phyxminiproblems-problemphyxmini0955-electricpotentialmagnitudequantity}
  \lean{PhyXMiniProblems.ProblemPhyXMini0955.ElectricPotentialMagnitudeQuantity}
  \uses{def:physics:phyx-mini-0955:phyxminiproblems-problemphyxmini0955-electricpotentialdimension}
  A nonnegative, unit-independent accelerating-potential magnitude
\end{definition}
\begin{proof}
  This is the declared model definition of Electric Potential Magnitude Quantity.
\end{proof}
\begin{definition}[Electric Field Magnitude Quantity]
  \label{def:physics:phyx-mini-0955:phyxminiproblems-problemphyxmini0955-electricfieldmagnitudequantity}
  \lean{PhyXMiniProblems.ProblemPhyXMini0955.ElectricFieldMagnitudeQuantity}
  \uses{def:physics:phyx-mini-0955:phyxminiproblems-problemphyxmini0955-electricfielddimension}
  A nonnegative, unit-independent electric-field magnitude
\end{definition}
\begin{proof}
  This is the declared model definition of Electric Field Magnitude Quantity.
\end{proof}
\begin{definition}[Electric Field Squared Magnitude Quantity]
  \label{def:physics:phyx-mini-0955:phyxminiproblems-problemphyxmini0955-electricfieldsquaredmagnitudequantity}
  \lean{PhyXMiniProblems.ProblemPhyXMini0955.ElectricFieldSquaredMagnitudeQuantity}
  \uses{def:physics:phyx-mini-0955:phyxminiproblems-problemphyxmini0955-electricfieldsquareddimension}
  A nonnegative, unit-independent squared electric-field magnitude
\end{definition}
\begin{proof}
  This is the declared model definition of Electric Field Squared Magnitude Quantity.
\end{proof}
\begin{definition}[Magnetic Flux Density Magnitude Quantity]
  \label{def:physics:phyx-mini-0955:phyxminiproblems-problemphyxmini0955-magneticfluxdensitymagnitudequantity}
  \lean{PhyXMiniProblems.ProblemPhyXMini0955.MagneticFluxDensityMagnitudeQuantity}
  \uses{def:physics:phyx-mini-0955:phyxminiproblems-problemphyxmini0955-magneticfluxdensitydimension}
  A nonnegative, unit-independent magnetic-flux-density magnitude
\end{definition}
\begin{proof}
  This is the declared model definition of Magnetic Flux Density Magnitude Quantity.
\end{proof}
\begin{definition}[nonnegative SIReadout]
  \label{def:physics:phyx-mini-0955:phyxminiproblems-problemphyxmini0955-nonnegativesireadout}
  \lean{PhyXMiniProblems.ProblemPhyXMini0955.nonnegativeSIReadout}
  Coherent-SI readout of a nonnegative dimensionful quantity
\end{definition}
\begin{proof}
  This is the declared model definition of nonnegative SIReadout.
\end{proof}
\begin{definition}[mass In Kilograms]
  \label{def:physics:phyx-mini-0955:phyxminiproblems-problemphyxmini0955-massinkilograms}
  \lean{PhyXMiniProblems.ProblemPhyXMini0955.massInKilograms}
  \uses{def:physics:phyx-mini-0955:phyxminiproblems-problemphyxmini0955-massquantity, def:physics:phyx-mini-0955:phyxminiproblems-problemphyxmini0955-nonnegativesireadout}
  Read a mass in coherent-SI kilograms
\end{definition}
\begin{proof}
  This is the declared model definition of mass In Kilograms.
\end{proof}
\begin{definition}[charge Magnitude In Coulombs]
  \label{def:physics:phyx-mini-0955:phyxminiproblems-problemphyxmini0955-chargemagnitudeincoulombs}
  \lean{PhyXMiniProblems.ProblemPhyXMini0955.chargeMagnitudeInCoulombs}
  \uses{def:physics:phyx-mini-0955:phyxminiproblems-problemphyxmini0955-chargemagnitudequantity, def:physics:phyx-mini-0955:phyxminiproblems-problemphyxmini0955-nonnegativesireadout}
  Read a charge magnitude in coherent-SI coulombs
\end{definition}
\begin{proof}
  This is the declared model definition of charge Magnitude In Coulombs.
\end{proof}
\begin{definition}[speed In Meters Per Second]
  \label{def:physics:phyx-mini-0955:phyxminiproblems-problemphyxmini0955-speedinmeterspersecond}
  \lean{PhyXMiniProblems.ProblemPhyXMini0955.speedInMetersPerSecond}
  \uses{def:physics:phyx-mini-0955:phyxminiproblems-problemphyxmini0955-speedmagnitudequantity, def:physics:phyx-mini-0955:phyxminiproblems-problemphyxmini0955-nonnegativesireadout}
  Read a speed magnitude in coherent-SI metres per second
\end{definition}
\begin{proof}
  This is the declared model definition of speed In Meters Per Second.
\end{proof}
\begin{definition}[electric Potential In Volts]
  \label{def:physics:phyx-mini-0955:phyxminiproblems-problemphyxmini0955-electricpotentialinvolts}
  \lean{PhyXMiniProblems.ProblemPhyXMini0955.electricPotentialInVolts}
  \uses{def:physics:phyx-mini-0955:phyxminiproblems-problemphyxmini0955-electricpotentialmagnitudequantity, def:physics:phyx-mini-0955:phyxminiproblems-problemphyxmini0955-nonnegativesireadout}
  Read an accelerating potential in coherent-SI volts
\end{definition}
\begin{proof}
  This is the declared model definition of electric Potential In Volts.
\end{proof}
\begin{definition}[electric Field In Volts Per Meter]
  \label{def:physics:phyx-mini-0955:phyxminiproblems-problemphyxmini0955-electricfieldinvoltspermeter}
  \lean{PhyXMiniProblems.ProblemPhyXMini0955.electricFieldInVoltsPerMeter}
  \uses{def:physics:phyx-mini-0955:phyxminiproblems-problemphyxmini0955-electricfieldmagnitudequantity, def:physics:phyx-mini-0955:phyxminiproblems-problemphyxmini0955-nonnegativesireadout}
  Read an electric-field magnitude in coherent-SI volts per metre
\end{definition}
\begin{proof}
  This is the declared model definition of electric Field In Volts Per Meter.
\end{proof}
\begin{definition}[electric Field Squared In Volts Squared Per Meter Squared]
  \label{def:physics:phyx-mini-0955:phyxminiproblems-problemphyxmini0955-electricfieldsquaredinvoltssquaredpermetersquared}
  \lean{PhyXMiniProblems.ProblemPhyXMini0955.electricFieldSquaredInVoltsSquaredPerMeterSquared}
  \uses{def:physics:phyx-mini-0955:phyxminiproblems-problemphyxmini0955-electricfieldsquaredmagnitudequantity, def:physics:phyx-mini-0955:phyxminiproblems-problemphyxmini0955-nonnegativesireadout}
  Read squared electric field in coherent-SI volts squared per metre squared
\end{definition}
\begin{proof}
  This is the declared model definition of electric Field Squared In Volts Squared Per Meter Squared.
\end{proof}
\begin{definition}[magnetic Flux Density In Teslas]
  \label{def:physics:phyx-mini-0955:phyxminiproblems-problemphyxmini0955-magneticfluxdensityinteslas}
  \lean{PhyXMiniProblems.ProblemPhyXMini0955.magneticFluxDensityInTeslas}
  \uses{def:physics:phyx-mini-0955:phyxminiproblems-problemphyxmini0955-magneticfluxdensitymagnitudequantity, def:physics:phyx-mini-0955:phyxminiproblems-problemphyxmini0955-nonnegativesireadout}
  Read magnetic flux density in coherent-SI teslas
\end{definition}
\begin{proof}
  This is the declared model definition of magnetic Flux Density In Teslas.
\end{proof}
\begin{definition}[Particle Species]
  \label{def:physics:phyx-mini-0955:phyxminiproblems-problemphyxmini0955-particlespecies}
  \lean{PhyXMiniProblems.ProblemPhyXMini0955.ParticleSpecies}
  Particle species distinguished by the scenario
\end{definition}
\begin{proof}
  This is the declared model definition of Particle Species.
\end{proof}
\begin{definition}[Charge Sign]
  \label{def:physics:phyx-mini-0955:phyxminiproblems-problemphyxmini0955-chargesign}
  \lean{PhyXMiniProblems.ProblemPhyXMini0955.ChargeSign}
  The sign of a particle's electric charge
\end{definition}
\begin{proof}
  This is the declared model definition of Charge Sign.
\end{proof}
\begin{definition}[Deflection Plate]
  \label{def:physics:phyx-mini-0955:phyxminiproblems-problemphyxmini0955-deflectionplate}
  \lean{PhyXMiniProblems.ProblemPhyXMini0955.DeflectionPlate}
  The two charged plates of the modern Thomson apparatus
\end{definition}
\begin{proof}
  This is the declared model definition of Deflection Plate.
\end{proof}
\begin{definition}[Graph Axis]
  \label{def:physics:phyx-mini-0955:phyxminiproblems-problemphyxmini0955-graphaxis}
  \lean{PhyXMiniProblems.ProblemPhyXMini0955.GraphAxis}
  Horizontal and vertical axes in the supplied graph
\end{definition}
\begin{proof}
  This is the declared model definition of Graph Axis.
\end{proof}
\begin{definition}[Plotted Quantity]
  \label{def:physics:phyx-mini-0955:phyxminiproblems-problemphyxmini0955-plottedquantity}
  \lean{PhyXMiniProblems.ProblemPhyXMini0955.PlottedQuantity}
  Physical quantities printed as the two graph-axis labels
\end{definition}
\begin{proof}
  This is the declared model definition of Plotted Quantity.
\end{proof}
\begin{definition}[Graph Anchor]
  \label{def:physics:phyx-mini-0955:phyxminiproblems-problemphyxmini0955-graphanchor}
  \lean{PhyXMiniProblems.ProblemPhyXMini0955.GraphAnchor}
  Four exact anchor points readable from the straight-line data plot
\end{definition}
\begin{proof}
  This is the declared model definition of Graph Anchor.
\end{proof}
\begin{definition}[ESquared Versus Voltage Graph]
  \label{def:physics:phyx-mini-0955:phyxminiproblems-problemphyxmini0955-esquaredversusvoltagegraph}
  \lean{PhyXMiniProblems.ProblemPhyXMini0955.ESquaredVersusVoltageGraph}
  \uses{def:physics:phyx-mini-0955:phyxminiproblems-problemphyxmini0955-graphaxis, def:physics:phyx-mini-0955:phyxminiproblems-problemphyxmini0955-plottedquantity, def:physics:phyx-mini-0955:phyxminiproblems-problemphyxmini0955-graphanchor}
  Literal presentation and coordinate information from image 955.png. Vertical coordinates are numbers in the printed scale 10⁸ V²/m²; they are not bare electric-field quantities
\end{definition}
\begin{proof}
  This is the declared model definition of ESquared Versus Voltage Graph.
\end{proof}
\begin{definition}[Muon Thomson Experiment]
  \label{def:physics:phyx-mini-0955:phyxminiproblems-problemphyxmini0955-muonthomsonexperiment}
  \lean{PhyXMiniProblems.ProblemPhyXMini0955.MuonThomsonExperiment}
  \uses{def:physics:phyx-mini-0955:phyxminiproblems-problemphyxmini0955-massquantity, def:physics:phyx-mini-0955:phyxminiproblems-problemphyxmini0955-chargemagnitudequantity, def:physics:phyx-mini-0955:phyxminiproblems-problemphyxmini0955-speedmagnitudequantity, def:physics:phyx-mini-0955:phyxminiproblems-problemphyxmini0955-electricpotentialmagnitudequantity, def:physics:phyx-mini-0955:phyxminiproblems-problemphyxmini0955-electricfieldmagnitudequantity, def:physics:phyx-mini-0955:phyxminiproblems-problemphyxmini0955-electricfieldsquaredmagnitudequantity, def:physics:phyx-mini-0955:phyxminiproblems-problemphyxmini0955-magneticfluxdensitymagnitudequantity, def:physics:phyx-mini-0955:phyxminiproblems-problemphyxmini0955-particlespecies, def:physics:phyx-mini-0955:phyxminiproblems-problemphyxmini0955-chargesign, def:physics:phyx-mini-0955:phyxminiproblems-problemphyxmini0955-deflectionplate, def:physics:phyx-mini-0955:phyxminiproblems-problemphyxmini0955-graphanchor, def:physics:phyx-mini-0955:phyxminiproblems-problemphyxmini0955-esquaredversusvoltagegraph}
  Independent physical data in the apparatus and at each plotted measurement. In particular, muonMass is not defined in terms of the answer or the graph
\end{definition}
\begin{proof}
  This is the declared model definition of Muon Thomson Experiment.
\end{proof}
\begin{definition}[Matches Scenario And Graph Readouts]
  \label{def:physics:phyx-mini-0955:phyxminiproblems-problemphyxmini0955-matchesscenarioandgraphreadouts}
  \lean{PhyXMiniProblems.ProblemPhyXMini0955.MatchesScenarioAndGraphReadouts}
  \uses{def:physics:phyx-mini-0955:phyxminiproblems-problemphyxmini0955-chargemagnitudeincoulombs, def:physics:phyx-mini-0955:phyxminiproblems-problemphyxmini0955-speedinmeterspersecond, def:physics:phyx-mini-0955:phyxminiproblems-problemphyxmini0955-electricpotentialinvolts, def:physics:phyx-mini-0955:phyxminiproblems-problemphyxmini0955-electricfieldsquaredinvoltssquaredpermetersquared, def:physics:phyx-mini-0955:phyxminiproblems-problemphyxmini0955-magneticfluxdensityinteslas, def:physics:phyx-mini-0955:phyxminiproblems-problemphyxmini0955-muonthomsonexperiment}
  Problem prose and literal content of the primary graph. The four dimensionful observables are tied to the scalar graph coordinates at the SI-readout boundary. This predicate contains no muon-mass value
\end{definition}
\begin{proof}
  This is the declared model definition of Matches Scenario And Graph Readouts.
\end{proof}
\begin{definition}[Has Elementary Charge Calibration]
  \label{def:physics:phyx-mini-0955:phyxminiproblems-problemphyxmini0955-haselementarychargecalibration}
  \lean{PhyXMiniProblems.ProblemPhyXMini0955.HasElementaryChargeCalibration}
  \uses{def:physics:phyx-mini-0955:phyxminiproblems-problemphyxmini0955-chargemagnitudeincoulombs, def:physics:phyx-mini-0955:phyxminiproblems-problemphyxmini0955-muonthomsonexperiment}
  Exact SI calibration of the elementary charge magnitude. This is independent external data, not a premise about the muon's requested mass
\end{definition}
\begin{proof}
  This is the declared model definition of Has Elementary Charge Calibration.
\end{proof}
\begin{definition}[Satisfies Muon Thomson Laws]
  \label{def:physics:phyx-mini-0955:phyxminiproblems-problemphyxmini0955-satisfiesmuonthomsonlaws}
  \lean{PhyXMiniProblems.ProblemPhyXMini0955.SatisfiesMuonThomsonLaws}
  \uses{def:physics:phyx-mini-0955:phyxminiproblems-problemphyxmini0955-massinkilograms, def:physics:phyx-mini-0955:phyxminiproblems-problemphyxmini0955-chargemagnitudeincoulombs, def:physics:phyx-mini-0955:phyxminiproblems-problemphyxmini0955-speedinmeterspersecond, def:physics:phyx-mini-0955:phyxminiproblems-problemphyxmini0955-electricpotentialinvolts, def:physics:phyx-mini-0955:phyxminiproblems-problemphyxmini0955-electricfieldinvoltspermeter, def:physics:phyx-mini-0955:phyxminiproblems-problemphyxmini0955-electricfieldsquaredinvoltssquaredpermetersquared, def:physics:phyx-mini-0955:phyxminiproblems-problemphyxmini0955-magneticfluxdensityinteslas, def:physics:phyx-mini-0955:phyxminiproblems-problemphyxmini0955-muonthomsonexperiment}
  The three modeling laws used in the calculation. Acceleration from the common initial speed is stated before substituting the zero-speed readout. The zero-deflection relation is the magnitude form of electric/magnetic Lorentz force balance. No field of this structure states the resulting mass formula
\end{definition}
\begin{proof}
  This is the declared model definition of Satisfies Muon Thomson Laws.
\end{proof}
\begin{definition}[Has Nondegenerate Muon Run]
  \label{def:physics:phyx-mini-0955:phyxminiproblems-problemphyxmini0955-hasnondegeneratemuonrun}
  \lean{PhyXMiniProblems.ProblemPhyXMini0955.HasNondegenerateMuonRun}
  \uses{def:physics:phyx-mini-0955:phyxminiproblems-problemphyxmini0955-massinkilograms, def:physics:phyx-mini-0955:phyxminiproblems-problemphyxmini0955-chargemagnitudeincoulombs, def:physics:phyx-mini-0955:phyxminiproblems-problemphyxmini0955-speedinmeterspersecond, def:physics:phyx-mini-0955:phyxminiproblems-problemphyxmini0955-electricpotentialinvolts, def:physics:phyx-mini-0955:phyxminiproblems-problemphyxmini0955-electricfieldinvoltspermeter, def:physics:phyx-mini-0955:phyxminiproblems-problemphyxmini0955-magneticfluxdensityinteslas, def:physics:phyx-mini-0955:phyxminiproblems-problemphyxmini0955-muonthomsonexperiment}
  Positivity and nondegeneracy conditions for the measured run
\end{definition}
\begin{proof}
  This is the declared model definition of Has Nondegenerate Muon Run.
\end{proof}
\begin{definition}[Answer Choice]
  \label{def:physics:phyx-mini-0955:phyxminiproblems-problemphyxmini0955-answerchoice}
  \lean{PhyXMiniProblems.ProblemPhyXMini0955.AnswerChoice}
  The four mass choices printed in the problem statement
\end{definition}
\begin{proof}
  This is the declared model definition of Answer Choice.
\end{proof}
\begin{definition}[answer Choice Mass In Kilograms]
  \label{def:physics:phyx-mini-0955:phyxminiproblems-problemphyxmini0955-answerchoicemassinkilograms}
  \lean{PhyXMiniProblems.ProblemPhyXMini0955.answerChoiceMassInKilograms}
  \uses{def:physics:phyx-mini-0955:phyxminiproblems-problemphyxmini0955-answerchoice}
  Numerical value in kilograms printed beside each answer choice
\end{definition}
\begin{proof}
  This is the declared model definition of answer Choice Mass In Kilograms.
\end{proof}
\begin{definition}[Agrees To Displayed Precision]
  \label{def:physics:phyx-mini-0955:phyxminiproblems-problemphyxmini0955-agreestodisplayedprecision}
  \lean{PhyXMiniProblems.ProblemPhyXMini0955.AgreesToDisplayedPrecision}
  Generic agreement with a displayed value within a positive half-unit
\end{definition}
\begin{proof}
  This is the declared model definition of Agrees To Displayed Precision.
\end{proof}
% --- Archon declaration topology end ---
% --- Archon physics formalization source end ---
